\begin{itemize}
\item \textbf{No real customer messages and no human gold labels.} Ottavo, its lexical stress set and Ottavo-R are synthetic and were
written and validated by Claude models; train and normal test come from the same generators. Scores near 0.99 on these sets are upper
bounds and show a ceiling effect. External validity rests on the two public benchmarks, which cover intent classification only.
\item \textbf{Descriptions written from training data, and a post-hoc zero-label comparison.} On Banking77 and MASSIVE-it the
descriptions (D) were written by agents that read labelled training examples, so D is not a zero-label condition there. The zero-label
budget comparison (label names only, family P2N) was added after the main analysis, following an external review; it was not
pre-declared, although it uses the same runs and procedure as the pre-declared family. On Ottavo D is the generation specification and
matches the data more closely than a real specification would.
\item \textbf{Possible contamination.} Banking77 and MASSIVE are public; Base Labs report training on public decision tasks. We cannot
rule out that the decision models (or the base models of every system) saw these data. The small effect of descriptions for Elephant-1
and Mammoth-1 on Banking77 is compatible with, but not evidence of, such exposure.
\item \textbf{Moving endpoints.} OpenAI Decisions was a public beta and the Baseten models were called through ``latest'' aliases. The
results describe the endpoints in October 2026; one re-run of OpenAI with label names differed by 0.012 macro-F1 across days.
\item \textbf{Test size.} Historical test samples have about 1{,}000 messages (601 and 312 for the Ottavo stress sets). Differences of
about one point are often not resolved; the full official tests (3{,}076 and 2{,}974 messages) are reported for the decision models and
the compact classifiers as a separate analysis.
\item \textbf{Recipes.} Hyperparameters were selected by cross-validation on seed 17 and frozen for the other seeds, with a small grid
(two learning rates, three durations). For mmBERT and Qwen3-1.7B at 5 examples per class the selected duration was the largest in the
grid. The full-data recipe (constant learning rate, epoch chosen on 10\% of the training data) proved suboptimal for Qwen3-0.6B on
Banking77, where the exploratory recipe scored 2.6 points higher (Section~\ref{sec:results}); the two recipes were not compared on the
other datasets or models. Decoder classifiers use a newly initialised linear head that does not read the category
descriptions; a variant that keeps the language-model head and the descriptions during fine-tuning was not tested, so conclusions about
fine-tuned decoders apply to this recipe only.
\item \textbf{Hardware and software.} Supervised GPU runs used a Colab Tesla T4 (fp16 mixed precision, Qwen3-8B at 4 bit, Qwen3-1.7B with
micro-batches of 4 and gradient accumulation to an effective batch of 16) and, for a few runs, an RTX 4050 laptop GPU. Library versions
were pinned only after the first Colab jobs (zero-shot scoring, SetFit, mmBERT) had run with the then-current releases. Latencies are not
comparable across API and local hardware.
\item \textbf{Generalist controls only on Ottavo.} Single-letter option scoring supports at most 26 options, so Qwen3-8B and Qwen3.5-9B were
evaluated only on the 24-class Ottavo task. The progenitor comparison does not isolate the effect of the decision training data from the
attention mask, the objective and the serving stack.
\item \textbf{Not tested.} Jev (waiting-list access only); self-hosting of Elephant-1 or Mammoth-1; distillation from a decision model
into a compact student; real-traffic costs, including the cost of labelling.
\end{itemize}
